\documentclass[a4paper,oneside,11pt]{article}

\usepackage[
    top=3cm,
    bottom=3cm,
    left=3cm,
    right=3cm,
    heightrounded,
    bindingoffset=5mm
]{geometry}

\pagestyle{plain}

\usepackage{amsfonts}
\usepackage{amsmath}
\usepackage{braket}
\usepackage{mathtools}
\usepackage{quantikz}
\usepackage{wrapfig}

\usepackage{hyperref}
\hypersetup{
    colorlinks,
    citecolor=black,
    filecolor=black,
    linkcolor=black,
    urlcolor=black
}

\setlength\parindent{0pt}

%\setlength\parindent{0pt}

\title{Topological Quantum Computing}
\author{Gergely Lendvay}
\date{\today}

\begin{document}
\maketitle\clearpage

\tableofcontents\clearpage

\newpage\section*{Introduction}
\subsection{Course Info}
\subsubsection{Teachers}
\noindent Konstantin Wernli -- kwernli@imada.sdu.dk

\noindent William Mistegard -- wem@imada.sdu.dk

\subsubsection{Exercise Classes}
Once a week with at least 3 solutions handed in.

\subsubsection{Exam}
\begin{enumerate}
    \item Assessment is based on four graded exercise sheets, and
    \item Final oral exam consisting of:
    \begin{itemize}
        \item 10--12 minutes presentation,
        \item 13--15 minutes discussion and questions,
        \item 4--5 possible topics, with 30 minutes of preparation time.
    \end{itemize}
    \item Tentative exam date: January 18, 2027
\end{enumerate}

\subsubsection{Books}
N/A

\newpage\section{First Lecture}
\subsection{Topological Quantum Computing}
Different fields:
\begin{itemize}
    \item Condensed matter physics: topological phases of matter
    \item Topological Quantum Field Theory
    \item TQFT's in math
    \item Anyon models/UMTC in math
\end{itemize}

\subsection{What is Topological Quantum Computing?}
Digital vs analog quantum computing.

\paragraph{Problem}
They are prone to errors and noise, such as spam error during state preparation, gate fidelities as well as decoherence.

The idea of TQC is to use very special physical systems where information is encoded in discrete global properties.
These are topological phases of matter.

Certain phases in $2+1$ dimensions, with 2 spatial dimensions and 1 time, admit anyons.
Anyons are a special class of (quasi) particles with statistics different from bosons and fermions.

In more than $2+1$ dimensions $T^2=1\rightarrow T=\pm1$
$+1=e^{i2\pi i}$ bosons
$-1=^{\pi i}$ fermions

In $2+1$ dimension $T_2\neq1$, meaning it can get any phase $T=e^{i\Theta}$

If $T=e^{i\Theta}\rightarrow$ abelian anyons, but there are non-abelian anyons where T can be a unitary acting on $\mathcal{H}$.
These are the unitary we want to use in TQC.

\subsection{Non-abelian anyons}
Described by $\ket{\psi}\in\mathcal{H}$ when exchanging anyons
\[\ket{\psi}\to U\ket{\psi},\]
for some unitary $U$.

Braids form a group.

Companies trying to develop non-abelian anyons: Microsoft and Nokia
Simulate anyons on quantum computer (topological codes): Google and Quantinuum

\subsection{Plan}
Math foundations:
\begin{itemize}
    \item Groups (braid),
    \item Representation,
    \item Category theory,
    \item Topology.
\end{itemize}

TQC:
\begin{itemize}
    \item Anyon models,
    \item Examples,
    \item Equivalence to digital quantum computing,
    \item Relation to TQFT's, etc.
\end{itemize}

TQC in practice:
\begin{itemize}
    \item Topological phases of matter,
    \item Experimental realizations.
\end{itemize}

Topological codes:
\begin{itemize}
    \item Surface codes,
    \item Other classes of topological codes.
\end{itemize}

\newpage\section{Second Lecture}
\subsection{}

\end{document}
